\section{Geometric norm maps for the nonordinary ray coefficient}
\label{sec:nonordinary-ray-norm}

This section continues AC-K33-ray-coefficient.  All coefficient quotients
are taken before the ordinary projector at \(p=269\).  The conclusion is
compatibility of the actual coefficient maps with explicit geometric
degeneracy maps.  It does not assert that a particular Rankin--Iwasawa
class survives these maps, or that its two singular regulators are
independent.

Fix \(K=\mathbb Q(\sqrt{-3})\), \(b\in\{1,2\}\), and
\(\rho_0=\operatorname{Ind}_{G_K}^{G_{\mathbb Q}}\chi^b\).
Let \(j\geq0\), and retain \(D_0\), its lattice, and the coefficient
ring \(\mathcal O\) from AC-K33-ray-coefficient.
Choose fixed tame integers \(M_0\mid N_0\), prime to \(p\), with
\(3M_0\mid N_0\), so that the corresponding modular curves are neat.
We consider squarefree ideals \(\mathfrak n\) prime to \(pN_0\);
both primes above a split rational prime are allowed to divide
\(\mathfrak n\).  Write
\[
 \Delta_{\mathfrak n}=\operatorname{Gal}(K(\mathfrak n)/K),\quad
 A_{\mathfrak n}=\mathcal O[\Delta_{\mathfrak n}],\quad
 \rho_{\mathfrak n}
   =\operatorname{Ind}_{G_K}^{G_{\mathbb Q}}
       (\chi^b\vartheta_{\mathfrak n}),
 \qquad
 L_{\mathfrak n}=N_0\operatorname N\mathfrak n\,p^a .
 \tag{RNN.1}
\]
Here \(K(\mathfrak n)\) is the maximal \(p\)-primary ray-class
extension of conductor dividing \(\mathfrak n\).
We use geometric Artin reciprocity: if
\(\mathfrak l\nmid\mathfrak n pN_0\), put
\[
 u_{\mathfrak l}
   =\chi^b(\operatorname{Frob}_{\mathfrak l}^{\rm geom})
      \vartheta_{\mathfrak n}
         (\operatorname{Frob}_{\mathfrak l}^{\rm geom})
       \in A_{\mathfrak n}^{\times}.
 \tag{RNN.2}
\]
The transition \(A_{\mathfrak n'}\to A_{\mathfrak n}\) is denoted
\(\operatorname{red}_{\mathfrak n'/\mathfrak n}\).
No character idempotent involving \(1/|\Delta_{\mathfrak n}|\) is used.

The precise source conventions are important.  For geometric
degeneracies we use KLZ, arXiv:1503.02888v6, Definitions 2.4.1--2.4.3,
the analytic parametrization in \S2.3, Remark 4.4.6, Theorem 4.5.1
and Remark 4.5.3.  For the standard \(R^1\pi_*\) bundle action we
use Emerton, \S2.4, printed pp.~12--13 of
\url{https://www.math.uchicago.edu/~emerton/pdffiles/coates.pdf}.
The Colmez--Wang normalization is \S1.2.3--1.2.4, printed pp.~21--23,
\S3.4.4, Lemma 3.9 and Proposition 3.11, and Remark 13.12 of
\url{https://webusers.imj-prg.fr/~pierre.colmez/KE.pdf}.
The local level assertions used below are also directly visible in
Schmidt, \emph{Some remarks on local newforms for GL(2)},
Proposition 2.1.2, Lemma 2.2.1, Theorem 2.3.5 and \S2.4,
especially (48)--(53), printed pp.~9--13 and 20--22:
\url{https://math.ou.edu/~rschmidt/papers/gl2.pdf}.

\begin{lemma}[AC-K33-ray-period-natural: compatible integral periods]
\label{AC-K33-ray-period-natural}
The periods in AC-K33-ray-period can be chosen naturally for all
quotients of finite abelian \(p\)-groups.  More precisely, there is
a single series \(b(T)\in W(\overline{\mathbb F}_p)[[T]]^\times\)
such that, for every such group ring \(A\) and every \(p\)-power-order
element \(z\in A^\times\), the convergent value
\[
 t_z=b(z-1),\qquad \varphi(t_z)=zt_z,\qquad t_1=1
 \tag{RNN.3}
\]
is a unit with integral inverse.  Its specialization under
\(A'\to A\) is \(t_{\bar z}\).
\end{lemma}
\begin{proof}
The map \(\varphi-1\) is onto on \(W(\overline{\mathbb F}_p)\).
Modulo \(p\) this is the surjectivity of \(x\mapsto x^p-x\) on the
algebraic closure.  If a solution has been found modulo \(p^r\),
adding \(p^r c\), with the residue of \(c\) solving the next
Artin--Schreier equation, lifts it modulo \(p^{r+1}\).  Completeness
gives a solution.  Choose \(b_0=1\), and recursively choose
\((\varphi-1)b_r=b_{r-1}\) for \(r\geq1\).  Then
\(b(T)=\sum b_rT^r\) satisfies
\(\varphi b(T)=(1+T)b(T)\).

If \(z^{p^s}=1\), then \(z-1\) is nilpotent modulo \(p\), hence
topologically nilpotent in the finite complete \(\mathcal O\)-algebra
\(A\).  Evaluation of \(b\) and its formal inverse therefore
converges, proves (RNN.3), and commutes with every coefficient
quotient.  Moreover \(\varphi^{p^s}t_z=t_z\), so its coefficients
belong to the unramified extension of degree \(p^s\).
All coefficients and their inverses are integral.  Thus both
operator norms are at most one, uniformly in \(s\).
\end{proof}

\begin{lemma}[AC-K33-ray-refinement: integral partial-depletion vectors]
\label{AC-K33-ray-refinement}
Let \(q\ne p\) be outside the fixed ramification set.
In the actual Emerton action on the \(\rho_{\mathfrak n}\)-part
of completed cohomology, the away-\(q\) representation is
\[
 \Pi_{\rm CW}(\rho_{\mathfrak n}^{\vee})^\sharp
       \simeq \Pi_{\rm CW}(\rho_{\mathfrak n}(1)),
 \qquad g^\sharp=(g^{-1})^t .
 \tag{RNN.4}
\]
In its upper Kirillov model put
\(D_q=\operatorname{diag}(q^{-1},1)\).
If both local characters are unramified with geometric Frobenius
values \(a,b\), there is a spherical vector \(v_{a,b}\) whose values
are
\[
 v_{a,b}(q^n u)=q^{-n}\sum_{i=0}^n a^{\,n-i}b^i
 \quad(n\geq0,\ u\in\mathbb Z_q^\times),
 \qquad v_{a,b}(q^n u)=0\quad(n<0).
 \tag{RNN.5}
\]
If the second character is unramified with value \(b\), and the
first is at most tamely ramified, there is an integral vector
\[
 v_b(q^n u)=(b/q)^n\quad(n\geq0),\qquad
 v_b(q^n u)=0\quad(n<0).
 \tag{RNN.6}
\]
The latter is fixed by \(K_1(q)\), including when the first
character specializes to an unramified character.  With
\(\delta_q=1_{\mathbb Z_q^\times}\), one has exact identities
\[
 (1-aq^{-1}D_q)v_{a,b}=v_b,\qquad
 (1-bq^{-1}D_q)v_b=\delta_q.
 \tag{RNN.7}
\]
For an unramified induction from the unramified quadratic extension,
whose squared Frobenius value is \(c\), these imply
\[
 (1-cq^{-2}D_q^2)v_{\rm sph}=\delta_q.
 \tag{RNN.8}
\]
All these vectors have integral values under every local group
translate and value one at the identity.  The compact vector is
fixed by \(K_1(q^2)\) for every tame ray specialization in question.
These assertions hold over the actual group ring, including its
conductor-drop fibers.
\end{lemma}
\begin{proof}
The inverse-transpose convention is literal in Colmez--Wang
Remark 13.12.  To check its effect on the norm, let \(\ell\) be a
\(\psi\)-Whittaker functional for \(\pi\), and put
\(w=\left(\begin{smallmatrix}0&1\\-1&0\end{smallmatrix}\right)\).
Since \(n(u)^\sharp=wn(u)w^{-1}\), the functional
\(\ell^\sharp=\ell\circ\pi(w^{-1})\) is a
\(\psi\)-Whittaker functional for \(\pi^\sharp\), and
\[
 W_v^\sharp(g)=W_v(w^{-1}g^\sharp).
 \tag{RNN.9}
\]
The map \(g\mapsto w^{-1}g^\sharp\) is a bijection of the group.
Thus (RNN.9) preserves the supremum over all Whittaker values
exactly.  No intertwining integral or epsilon-factor is divided.

The normalized principal-series parameters of
\(\Pi_{\rm CW}(\rho^\vee)^\sharp\) are those of
\({\rm LL}(\rho)|\det|^{1/2}\).  Equivalently this is the
Colmez--Wang model of \(\rho(1)\); geometric Frobenius on
\(\rho(1)\) has values \(a/q,b/q\).  This proves (RNN.4) and (RNN.5).
It also checks the Hecke normalization independently: the raw
double-coset sum at the identity is \(a+b\), whereas the
Colmez--Wang operator in \S1.2.3 contains the displayed factor
\(1/q\).

We give an integral proof of (RNN.6), because characterwise
diagonalization would not prove its uniform bound.
In the notation of Colmez--Wang \S3.4.4, use the induced
presentation \(I(\xi_2,\xi_1)\), with \(\xi_2\) unramified,
\(\xi_2(q)=b/q\), and both characters unit-valued for the
\(p\)-adic norm.  In its affine chart take
\(\bar\phi=\xi_1(-1)^{-1}1_{\mathbb Z_q}\).
The change of variables in the complete proof of Proposition 3.11
gives
\[
 K\bar\phi(y)
   =\xi_1(-1)\xi_2(y)
       \int_{\mathbb Q_q}\bar\phi(u)\psi(yu)\,du
   =1_{\mathbb Z_q}(y)\xi_2(y).
 \tag{RNN.10}
\]
Here \(\operatorname{vol}(\mathbb Z_q)=1\) and the Gauss factor
of the unramified \(\xi_2\) is one.
The induced section has integral values on every group element:
its chart is integral and all inducing characters and powers of
\(q\) are \(p\)-adic units.  The same is true after every translate.
Each truncated Jacquet integral is consequently integral.
For a fixed translate its characteristic-zero Jacquet integrals
stabilize on every generic component.  There are finitely many
components, so a common truncation works; injectivity of the
group ring into its rational components proves stabilization
over the group ring itself.  This proves the all-translate bound.
Alternatively the formula for the fractional linear action in
Lemma 3.9 proves the same assertion directly.

That formula also proves \(K_1(q)\)-invariance of this section.
For \(c\in q\mathbb Z_q\), \(d\equiv1\bmod q\), the fractional
linear transformation preserves \(\mathbb Z_q\) and its
complement; on \(\mathbb Z_q\), its character factor has
\(\det(g)/(a-cx)\equiv d\equiv1\bmod q\).
The unramified character is trivial on units.
The proof includes the unramified specialization of \(\xi_1\).

Equation (RNN.7) follows at each \(q^n u\) by subtracting successive
complete homogeneous polynomials, and holds also at \(n=0\)
because \(D_qv\) is zero there.  It is a polynomial identity,
so equal roots cause no problem.  In the inert case the trace
is zero, the squared root is \(c\), and (RNN.8) follows without
adjoining a square root.

The compact vector and its all-translate integral norm follow
from AC-K33-ray-Schwartz, followed by the unit-valued unramified
twist \(|\det|_q\) and (RNN.9).
For its level, each ray character at \(q\ne p\) is tame.
For split primes, the two conductor exponents are at most one.
If both are ramified, Schmidt's Proposition 2.1.2 and the
calculation (49) in \S2.4 identify the compact function with
the newvector at level at most \(q^2\).
If one or both are unramified, (RNN.7) and the elementary inclusion
\(D_q^{-1} K_1(q^2)D_q\subset K_1(q)\) give the same bound.
For an inert prime, a nontrivial tame character of the quadratic
extension gives either a dihedral supercuspidal of conductor two,
or two tamely ramified characters of conductor one.  Schmidt's
Theorem 2.3.5 and its \S2.4 Kirillov computation give the compact
newvector in the first case; the preceding principal-series
argument gives the second.  If the character becomes unramified,
(RNN.8) gives the bound.
All these identities can be checked in the rational components
and then in the group ring by injectivity.  Their integral norm
has already been established before this check.
\end{proof}

\begin{lemma}[AC-K33-ray-degeneracy: the exact coefficient scalar]
\label{AC-K33-ray-degeneracy}
Put \(W=R^1\pi_*\mathcal O\), and use the literal KLZ
degeneracy \(\operatorname{pr}_2\) and its canonical degree-\(q\)
isogeny \(\lambda\).
On compact-support cohomology with coefficients
\(\operatorname{Sym}^jW\), the transpose of the KLZ pushforward
on \(\operatorname{TSym}^j\mathcal H(1)\) is
\[
 (\operatorname{pr}_{2,*})^t
       =q^j\,D(q),\qquad
 D(q)=\operatorname{diag}(q^{-1},1)\in{\rm GL}_2(\mathbb A_f).
 \tag{RNN.11}
\]
Here \(D(q)\) is the finite rational diagonal, with its unit
components away \(q\).  The identity holds for finite
\(\mathcal O/p^r\) coefficients at every level before completion.
On a tensor whose non-\(q\) components are fixed by these units,
\(D(q)\) can be replaced by its \(q\)-local component \(D_q\).
For the degree-\(q^2\) degeneracy \(\operatorname{pr}_3\), the
corresponding identity is \(q^{2j}D(q)^2\).
\end{lemma}
\begin{proof}
Use the analytic parametrization of KLZ \S2.3:
\((a,\tau)\) represents
\((\mathbb C/(\mathbb Z+\mathbb Z\tau),a\tau/M,1/N)\).
The second degeneracy is
\((a,\tau)\mapsto(a,q\tau)\), with
\(\lambda:z\mapsto qz\).
Order the homology cycles as \((\tau,1)\), and use their dual
basis for \(W\).  The maps on the two homology lattices are
\[
 \lambda_*=\begin{pmatrix}1&0\\0&q\end{pmatrix},
 \qquad
 \lambda'_*=\begin{pmatrix}q&0\\0&1\end{pmatrix},
 \qquad \lambda'\lambda=[q].
 \tag{RNN.12}
\]
The same diagonal matrices give \(\lambda^*\) and \(\lambda'^*\)
in the corresponding dual bases, in their respective directions.

Emerton's bundle is
\[
 {\rm GL}_2(\mathbb Q)\backslash
  \bigl((\mathbb C\setminus\mathbb R)\times
       {\rm GL}_2(\mathbb A_f)\times\operatorname{Sym}^j\mathcal O^2\bigr)/K,
\]
where right translation by \(h\) sends the coefficient to
\(h_p^{-1}v\), and the right \(K\)-quotient uses \(k_p^{-1}v\).
For \(h=D(q)\), move the rational diagonal
\(\operatorname{diag}(q,1)\) to the left.  The base point becomes
\(q\tau\), and the forward coefficient map on \(W\) is
\(\operatorname{diag}(q,1)=\lambda'^*\).
The induced pullback on cohomology uses the inverse of that
forward bundle map.  Since
\(\lambda^*\lambda'^*=q\), it is
\(q^{-1}\lambda^*\) on \(W\), hence \(q^{-j}\lambda^*\)
on \(\operatorname{Sym}^jW\).
Poincare adjunction for the actual KLZ map
\(\operatorname{TSym}^j\lambda_*\), followed by curve trace,
has transpose precisely the curve pullback followed by
\(\operatorname{Sym}^j\lambda^*\).
This proves (RNN.11).

This is a calculation of the coefficient local-system morphisms,
not just of their action on holomorphic forms.
The charts run over all Weil-pairing components and the formulas
respect the level identifications.  Equivalently (RNN.12) expresses
the same comparison using the algebraic isogeny and its dual.
Finite-coefficient Betti--etale comparison identifies these maps
on the actual etale local systems.  Both correspondences are
defined on the full \(\mathbb Q\)-scheme, so the identity respects
its Galois action.  All matrices and their inverses are integral
at \(p\), since \(q\ne p\).
It therefore holds on finite coefficient complexes, on
compact-support cohomology, and after either limit.
No assertion that classical vectors are dense in an Artin
completed-cohomology quotient is used; no parabolic or boundary
comparison is needed for (RNN.11).
Composition proves the degree-\(q^2\) formula.
\end{proof}

\begin{theorem}[AC-K33-ray-norm: compatible actual ray coefficient maps]
\label{AC-K33-ray-norm}
The quotients of AC-K33-ray-coefficient can be chosen on the
levels \(Y(M_0,L_{\mathfrak n}p^r)\) of (RNN.1), with one
\(a\) and one denominator bound \(C\), so that they are compatible
with the following geometric norm maps.
Write
\[
 \mathcal I_{\mathfrak n,j}
 =\varprojlim_r H^1_{\rm et}
   \bigl(Y(M_0,L_{\mathfrak n}p^r)_{\overline{\mathbb Q}},
               \operatorname{TSym}^j\mathcal H(1)\bigr),
 \qquad
 T_{\mathfrak n}
  =\rho_{\mathfrak n}^{\vee_{A_{\mathfrak n}}}
                         \otimes_{\mathcal O}D_0^\vee.
 \tag{RNN.13}
\]
There are bounded equivariant surjections
\[
 Q_{\mathfrak n}:\mathcal I_{\mathfrak n,j}[1/p]
              \longtwoheadrightarrow T_{\mathfrak n}[1/p],
 \qquad
 Q_{\mathfrak n}(\mathcal I_{\mathfrak n,j})
       \subset p^{-C}T_{\mathfrak n}^{\circ},
 \tag{RNN.14}
\]
where
\(T_{\mathfrak n}^{\circ}
=\rho_{\mathfrak n}^{\vee_{A_{\mathfrak n}}}
 \otimes_{\mathcal O}(D_0^\circ)^\vee\),
with \(U_p'=0\) and inverse diamond \(u^{-j-1}\), such that:
if \(\mathfrak n'=\mathfrak n\mathfrak l\) and
\(\mathfrak l\nmid\mathfrak n\), then
\[
 \mathcal N_{\mathfrak n}^{\mathfrak n'}
 =\begin{cases}
   \operatorname{pr}_{1,*}
      -u_{\mathfrak l}q^{-j-1}\operatorname{pr}_{2,*},
     &q=\mathfrak l\bar{\mathfrak l}\text{ split},\\[2mm]
   \operatorname{pr}_{1,*}
      -u_{\mathfrak l}q^{-2j-2}\operatorname{pr}_{3,*},
     &\mathfrak l=(q)\text{ inert},
  \end{cases}
 \tag{RNN.15}
\]
and
\[
 Q_{\mathfrak n}\mathcal N_{\mathfrak n}^{\mathfrak n'}
    =\operatorname{red}_{\mathfrak n'/\mathfrak n}
           Q_{\mathfrak n'}.
 \tag{RNN.16}
\]
In (RNN.15), coefficients act in \(A_{\mathfrak n}\), after
specialization of the source Hecke algebra.  Thus (RNN.16) is
an equality of maps from the higher actual cohomology tower
to \(T_{\mathfrak n}[1/p]\).
The maps for different added primes commute.  Raising an
already present prime power uses \(\operatorname{pr}_{1,*}\);
no new \(p\)-primary tame ray character is introduced.

These are geometric norm maps for the actual second-factor
coefficient.  Tensoring with the fixed first-factor projection
gives the rank-four coefficient over \(\mathbb Q\), whose Shapiro
coefficient over \(K\) has rank two.
The result does not identify the Euler polynomial of a chosen
Rankin class by itself.
\end{theorem}
\begin{proof}
\emph{1. Local vectors at the prescribed levels.}
At a split prime \(q=\mathfrak l\bar{\mathfrak l}\), use the
spherical vector when neither prime divides \(\mathfrak n\);
use \(v_b\) when only \(\mathfrak l\) divides \(\mathfrak n\),
where \(b\) is the Frobenius value of the still unramified
\(\bar{\mathfrak l}\)-character; and use \(\delta_q\) when
both divide \(\mathfrak n\).
At an inert prime use the spherical vector when it does not
divide \(\mathfrak n\), and \(\delta_q\) otherwise.
The local levels are respectively \(1,q,q^2\) at a split
prime and \(1,q^2\) at an inert prime.  They divide
\(\operatorname N\mathfrak n\).
At the fixed tame prime \(3\), use the same unramified-tail
construction in the local representation consisting of one
unramified character and its ramified quadratic twist.
Its level is \(3\), independent of \(\mathfrak n\).
At any other fixed auxiliary level prime use the spherical
vector.  All upper diagonal unit matrices fix these profiles.

Use at \(p\) the simultaneous vectors \(\Phi_{j,d}\) of
AC-K33, \(d\in D_0\).  Choose their unramified periods by
(RNN.3), including the natural square root of the \(p\)-power
Frobenius element as in (RCQ.10).
They are \(P_0\)-fixed and killed by \(U_p\), at one fixed
depth.  Thus the unit components away \(q\) in \(D(q)\)
fix the entire tensor.

The integral norms of the away vectors are at most one,
uniformly in every ray group.  Their value one gives the
same lower bound used in AC-K33-ray-coefficient.
Formula (RNN.9) transports these precise bounds to the
Colmez--Wang convention before applying its global integral
saturation (3.30).
The uniform upper and lower bounds at \(p\) are unchanged.
Hence the construction of that theorem, followed by bounded
duality, gives (RNN.14) with one constant \(C\).

These vectors lie at the stated \(\Gamma_1\)-type auxiliary
levels, rather than only at unspecified principal levels.
To justify the integral descent, pass temporarily to a neat
principal tame cover on which they are defined.
It is a finite etale torsor over the stated level.
For the finite coefficient sheaves on the affine curve,
\(H_c^0=0\).  The degree-one compact-support Cartan--Leray
edge is consequently an isomorphism onto the invariants.
It uses no group-order denominator.
The same finite-level argument and completion used in
AC-K33 descend the fixed vectors integrally.
The residual Artin ideal is the same non-Eisenstein ideal
throughout; the compact-support to ordinary comparison used
in AC-K33 therefore applies with its unchanged integral
normalization.

\emph{2. Coefficient reduction before character splitting.}
For a ray quotient \(A'\twoheadrightarrow A\), duality gives
the isometric inclusion
\[
 A^\dagger=\operatorname{Hom}_{\mathcal O}(A,\mathcal O)
           \hookrightarrow
 (A')^\dagger.
 \tag{RNN.17}
\]
In group-element dual bases this repeats the coefficients
on each fiber of the quotient.  Equivalently, after the
identity-coefficient Frobenius pairing \(A\simeq A^\dagger\),
it multiplies a lift by the sum of the kernel elements.
This is an integral isometry, not a character idempotent.
For \(h\in A'\) and \(v\in A^\dagger\),
\(h\,i(v)=i(\bar h\,v)\).
The induced Galois bases and the periods (RNN.3) commute
with this inclusion.  All local profile identities may
therefore be specialized in the actual coefficient dual.

We spell out why the construction is compatible with the
actual completed-cohomology occurrence, not just with an
abstract product of local representations.
For any finite collection of ray moduli, take one common
tame level and its actual Hecke algebra \(T\).
Each \(T\to A_{\mathfrak n}\) is the trace-surjective map
of (RCQ.6)--(RCQ.7), and its dual embeds into the same
\(T^\dagger\).
The rational finite Artin fibers have no exceptional
minimum representation, so Colmez--Wang Theorem 3.34(i)
identifies their rational spaces with the full local
representations.  The all-translate bounds put the
vectors in the literal integral intersection (3.30).

Colmez--Wang's actual cohomology map is uniquely
\(\iota_T=(K_H^G)^{-1}K_{\rm Aut}^G\)
(Theorem 13.11 and its proof).  In \S9.3.2 its
Hecke-enrichment is explicitly
\(K_H^T(v)(x)(t)=K_H(tv)(x)\).
Consequently the inclusions (RNN.17) and the corresponding
coefficient reductions commute with both sides of this
identity.  The literal products of local Whittaker
functions agree after specialization; no rational
averaging is performed.
The injectivity of \(K_H^G\), with its unchanged
non-Eisenstein residual hypothesis, then proves equality
of the actual cohomology vectors.
This also proves independence of the chosen common
Hecke level: compare two such choices at a third level.
The scalar ambiguity in a Galois realization cancels
against its dual local factor, as in Colmez--Wang
Remark 13.2(i).  Thus these comparisons glue for all
finite ray diagrams; no compactness of an unspecified
space of continuous maps is being assumed.

\emph{3. The geometric maps are the local depletion maps.}
Let first \(\mathfrak l\) be split, and put
\(a=u_{\mathfrak l}\in A_{\mathfrak n}\).
After reduction of the higher coefficient, its chosen
local profile is \(v_b\) if the lower profile was
\(v_{a,b}\), and is \(\delta_q\) if the lower profile was
\(v_a\).  In the second case the notation simply
interchanges the roles of the two primes.
In both cases (RNN.7) says that the higher profile is
\((1-aq^{-1}D_q)\) times the lower profile.
On the actual dual tower, (RNN.11) gives
\[
 (\mathcal N_{\mathfrak n}^{\mathfrak n'})^t
       =1-aq^{-j-1}q^jD(q)
       =1-aq^{-1}D_q
 \tag{RNN.18}
\]
on our tensors.  This is exactly the same identity,
including the geometric Frobenius orientation.

For an inert \(\mathfrak l=(q)\), lower Frobenius has trace
zero and squared value \(u_{\mathfrak l}\).
Equations (RNN.8) and (RNN.11) give instead
\[
 (\mathcal N_{\mathfrak n}^{\mathfrak n'})^t
      =1-u_{\mathfrak l}q^{-2j-2}q^{2j}D(q)^2
      =1-u_{\mathfrak l}q^{-2}D_q^2 .
 \tag{RNN.19}
\]
Step 2 promotes these identities of the actual integral
Whittaker vectors to identities in compact cohomology.
Taking the bounded dual proves (RNN.16).
The dual of the target reduction is precisely (RNN.17);
the multiplication by \(u_{\mathfrak l}\) has the same
orientation in the Frobenius pairing, not its inverse.

Different rational-prime correspondences commute on the
moduli problem.  For the two primes above the same split
\(q\), the two orders give the identical polynomial
\((1-aq^{-1}D_q)(1-bq^{-1}D_q)\).
The canonical cyclic quotient diagrams for
\(\operatorname{pr}_1,\operatorname{pr}_2\) commute as
well, so this is equality of the actual composed maps.
Finally, at \(q\ne p\) the higher unit groups are pro-\(q\).
They have no nontrivial \(p\)-primary ray character.
An already present prime therefore leaves the coefficient
and profile unchanged, and ordinary pullback is the
transpose of \(\operatorname{pr}_{1,*}\).
This proves the stated prime-power assertion.
\end{proof}

\begin{remark}[The remaining class and regulator questions]
Theorem AC-K33-ray-norm is a statement about actual geometric
coefficient maps, with a fixed integral denominator.
To obtain the desired \(K\)-Euler system from the actual
preordinary Rankin--Iwasawa classes, one must still insert
these maps into their precise asymmetric norm identities
and verify the resulting rank-two Euler polynomial with
all torsor, diamond and cyclotomic factors.
One must then prove that the resulting classes have
nonzero, independent singular images at \(269\).
Neither conclusion follows from the existence or
surjectivity of the coefficient maps, and taking the
wedge of two scalar norm systems would not supply the
required rank-two norm identity.
\end{remark}
